% Generated draft: requires mathematical and source review.
\documentclass[11pt]{article}
\usepackage[margin=25mm]{geometry}
\usepackage{fontspec}
\setmainfont{Latin Modern Roman}
\usepackage{amsmath,amssymb,mathtools,unicode-math}
\setmathfont{Latin Modern Math}
\usepackage{enumitem,needspace,xurl,hyperref}
\setlength{\parindent}{0pt}
\setlength{\parskip}{5pt}
\setlength{\emergencystretch}{3em}
\newcommand{\N}{\mathbb{N}}
\newcommand{\Z}{\mathbb{Z}}
\newcommand{\Q}{\mathbb{Q}}
\newcommand{\R}{\mathbb{R}}
\newcommand{\C}{\mathbb{C}}
\newcommand{\F}{\mathbb{F}}
\newcommand{\A}{\mathbb{A}}
\newcommand{\PP}{\mathbb{P}}
\newcommand{\E}{\mathbb{E}}
\newcommand{\eps}{\varepsilon}
\newcommand{\dd}{\,\mathrm{d}}
\begin{document}
\section*{UnsolvedMath ID: 30004144}
\textbf{Source problem code:} catalogue entry
\subsection{Definitions and mathematical statement}
Let \(X\) be a nonsingular projective real algebraic curve and let \(A\) be a simple real divisor on \(X\), as in the supplied record. Write \(X_{\C}=X\times_{\R}\C\), and let \(\sigma\) denote complex conjugation on \(X_{\C}(\C)\). A divisor on \(X_{\C}\) is a finite formal sum \(D=\sum_p n_p[p]\), with \(n_p\in\Z\). Conjugation invariance means \(\sigma_*D=D\); this permits conjugate pairs of nonreal points and is distinct from requiring every support point to belong to \(X(\R)\).

The record does not define ``simple real divisor.'' If it means a reduced effective conjugation-invariant divisor, then \(A\) has the form \(A=[p_1]+\cdots+[p_d]\) with distinct points and conjugation-invariant support. This is a conditional interpretation, not an additional supplied hypothesis. The record does not specify whether nonreal conjugate pairs are allowed or whether effectivity is required.

The requested motion is within the family
\[
\mathcal M_A=\{B:B\text{ is a simple real divisor on }X,\ B\sim A\},
\]
where \(\sim\) denotes linear equivalence. For divisors on a specified curve over a specified field, linear equivalence means that their difference is the divisor of a nonzero rational function. The record does not specify whether the relevant linear system and rational functions are taken over \(\R\) or after extension to \(\C\).

\textbf{Problem.} For each such pair \((X,A)\), describe the motions of \(A\) through simple real divisors linearly equivalent to \(A\), determine
\[
\pi_1\bigl(\operatorname{Emb}(A,X),a_0\bigr),
\]
and characterize the possible images of the stated trace homomorphism
\[
\operatorname{tr}:\pi_1\bigl(\operatorname{Emb}(A,X),a_0\bigr)\longrightarrow H_1(X).
\]
Here \(a_0\) must represent the initial configuration \(A\). The source does not define \(\operatorname{Emb}(A,X)\), its topology, or its relationship to \(\mathcal M_A\). In particular, it does not specify whether configurations are labelled or unlabelled, or whether the space is restricted to divisors linearly equivalent to \(A\). Thus the displayed fundamental group is a target requiring these definitions, rather than a fully specified invariant in the supplied record.

Under an unlabelled reduced-effective-divisor interpretation, the usual candidate for a trace sends a loop of configurations to the homology class of the sum of the trajectories of its points; the sum is a cycle even if the points are permuted. This construction still requires a specified ambient topological space and homology coefficients. The notation \(H_1(X)\) in the record does not distinguish \(H_1(X(\R);\Z)\) from \(H_1(X_{\C}(\C);\Z)\), or specify other coefficients. No choice between these targets is asserted here.
\subsection{Short English statement}
The problem asks how a simple real divisor can move while staying simple, real, and linearly equivalent to its initial value, and asks for the fundamental group of the resulting embedding space and the image of its trace map. The meanings of simplicity, the embedding space, and the homology target require clarification before this becomes a fully specified mathematical problem. The supplied status is \emph{partially solved}; this is a curation label, not an independently verified conclusion.
\subsection{Sources}
\begin{itemize}
\item Local source registry: \texttt{data/unsolvedmath/problems.json}, with \texttt{catalog/problems.json} as the additional-record registry; permanent UnsolvedMath ID \(30004144\), title \emph{Motion and Trace of Real Divisors}, source problem code \texttt{catalogue entry}. This draft uses the supplied complete registry record.
\item Supplied report citation: \emph{Tropical Geometry: new directions} (2020), Oberwolfach Reports, pp.\ 1245--1308, DOI \texttt{10.4171/owr/2019/21}. \url{https://doi.org/10.4171/owr/2019/21}. The record separately gives the proposed year as 2019 and supplies no named proposer.
\item Supplied literature source, labelled ``Tropical geometry problem report'': \url{https://ems.press/content/serial-article-files/46800}.
\item Literature assessment checked at \texttt{2026-08-22}: supplied dated curation, not independently verified. It labels the problem partially solved and reports that a real-divisor trace theorem determines important cases, including trivial and infinite cyclic trace images, while the full fundamental group of the embedding space remains undetermined in general. The supplied background explicitly requests independent verification of time-sensitive status and cited claims.
\item The registry labels the statement \texttt{corrected\_verified} and reports a manual semantic, typographic, and LaTeX review. That supplied review label does not resolve the missing definitions identified in this draft. The report links have not been independently consulted here.
\end{itemize}
\end{document}
